\section*{S31. Additional projection development and evaluation limits}
This section records completed development observations while confirmation and
calibration remain in progress. The development matrix comprises 32 previously
studied objects per interface, four conditions, and three sampling seeds
(42, 2026, 31415), for 384 object--condition--seed cells per interface. The two
interfaces use different object sets and are not pooled as paired subjects.
Sampling seeds do not estimate training-checkpoint variation.

For matched noisy latents, let $f_u,f_c$ denote negative/positive conditional
predictions with geometry and $b_u,b_c$ the corresponding predictions without
geometry. Define $a=b_c-b_u$, $g_0=f_u-b_u$, and
$h=(f_c-f_u)-a$. The native guided prediction is $b_u+g_0+s(a+h)$.
The development interaction candidate removes the negatively parallel component
of $h$ along $a$ and retains $g_0$ algebraically. Total-geometry grid and
view projections are separate comparison conditions. These are operations in
noise-prediction space; $a$ is not a pure appearance-loss gradient. Algebraic
decomposition and vector projection alone do not establish method novelty.

On the unmodified first-interface seed-42 trajectories, 32 objects provide
50 observations each. The object-averaged cosine between the guided geometry
contribution and $a$ is $-0.6426$ (95\% object-bootstrap interval
$[-0.6611,-0.6235]$); all 1600 measured whole-grid rows are negative.
The auxiliary measurement preserves final native pixels exactly. This validates
measurement without disturbing the trajectory; negative cosine does not prove
that the geometry contribution harms final appearance. Two MV-Adapter objects
also show negative whole-object cosines at all 50 steps, with respective means
$-0.4965$ and $-0.4984$. Two objects do not establish population generalization.

Table~\ref{tab:development012} reports projection-minus-native differences on
first-interface development objects. Seeds and views are nested within objects;
intervals use 20,000 object-bootstrap resamples and are descriptive rather than
multiplicity-adjusted confirmatory tests. Source agreement improves, but the
nonreference LPIPS intervals cross zero. No noninferiority margin was specified,
and absence of a detected difference is not evidence of no quality loss.

\begin{table}[ht]
\centering\scriptsize
\begin{tabular}{lccc}
\toprule
Condition & Source FG-PSNR (dB) & Source $\Delta E_{00}$ & Nonreference FG-LPIPS \\
\midrule
Total geometry, grid & $+2.74$ $[+1.57,+3.97]$ & $-1.09$ $[-2.09,-0.38]$ & $-0.0016$ $[-0.0044,+0.0012]$ \\
Total geometry, view & $+3.01$ $[+1.80,+4.30]$ & $-1.16$ $[-2.16,-0.46]$ & $+0.0014$ $[-0.0017,+0.0047]$ \\
Interaction only, grid & $+2.79$ $[+1.56,+4.09]$ & $-1.06$ $[-2.07,-0.34]$ & $-0.0014$ $[-0.0046,+0.0017]$ \\
\bottomrule
\end{tabular}
\caption{Exploratory development differences and 95\% object intervals.
PSNR is higher-better; color difference and LPIPS are lower-better. All three
sampling seeds are averaged within each of the 32 objects. These are not new
independent confirmation objects.}
\label{tab:development012}
\end{table}

Second-interface endpoint interpretation needs additional calibration. Native
view 0 has azimuth $-90^\circ$ and view 1 has azimuth $0^\circ$; source-view roles
cannot be transferred from the first interface by index alone. In the current
32-object set, eight target assets use a white texture fallback. The source-alpha
and candidate native-view-1 target masks have median pixel IoU 0.2606, so azimuth
agreement alone does not establish pixel registration. Earlier full-frame
reference comparison is also confounded by background: reference-alpha
background contributes a median 98.2\% of full-frame squared error for native
seed-42 view-1 outputs. A separate foreground/background-controlled diagnostic
has completed, but its pose-registration limitation prevents interpreting it as
validated reference fidelity or changing the confirmation endpoints.

\paragraph{Completed source-view confirmation and remaining tradeoffs.}
The 64-object, four-condition, three-seed matrix is complete (768 outputs).
The cohort excludes the 32 development objects but remains within historically
exposed Fresh C. The pre-generation specified view-projection versus native
comparisons give source FG-PSNR $+4.283$ dB (object-bootstrap 95\% CI
$[+3.409,+5.205]$; 61/64 object means favorable) and CIEDE2000 $-2.230$
($[-3.061,-1.563]$; 57/64 favorable). Both two-sided paired Monte Carlo
sign-flip tests have Holm-adjusted $p\approx 0.000010$ across the two primaries.
Seeds and views are nested within objects, not treated as independent samples.

Descriptive nonreference-view differences are FG-PSNR $-0.282$ dB
(pointwise 95\% CI $[-0.672,+0.102]$), FG-LPIPS $+0.000269$
($[-0.001533,+0.002073]$), CIEDE2000 $+0.425$
($[-0.157,+1.042]$), and FG-SSIM $+0.01062$
($[+0.00171,+0.01957]$). These secondary intervals are unadjusted.
Non-detection does not establish non-inferiority: no margin was specified.
The result confirms source-view gains for this comparison, not universal
multi-view gains, superiority of interaction projection, APG superiority,
or cross-interface transfer. APG's enabled norm-clipping path initially
failed because a matrix-norm operation received three reduction dimensions.
An explicit vector norm fixes that implementation error without changing the
frozen parameters; failed attempts are retained and a separate attempt is running.
